% =============================================================================
%  Lab 4 - Datapath Support Components
% =============================================================================

\chapter{Lab 4 --- Datapath Support Components}

\section{Objective}
% Two or three sentences, your own words.

\section{Specification}

\begin{table}[htbp]
\centering
\caption{Compressed table (Type C) for \texttt{ImmedXtnd}.
$x^{n}$ means bit $x$ repeated $n$ times.}
\label{tab:lab04-truth}
\begin{tabular}{ccc||ll}
\toprule
\# & \texttt{Xsel} & \texttt{Xlsft} & \texttt{Xtnd32} & used by \\
\midrule
0 & 0 & 0 & $0^{16}\,\&\,I$ & \texttt{LUI} \\
1 & 1 & 0 & $I_{15}^{16}\,\&\,I$ & \texttt{ADDI}, \texttt{LW}, \texttt{SW} \\
2 & 1 & 1 & $I_{15}^{14}\,\&\,I\,\&\,\texttt{"00"}$ & \texttt{BLEZ} \\
3 & 0 & 1 & $0^{14}\,\&\,I\,\&\,\texttt{"00"}$ & unused, but defined \\
\bottomrule
\end{tabular}
\end{table}

\begin{table}[htbp]
\centering
\caption{\texttt{std\_logic} resolution for two drivers on one wire.}
\label{tab:lab04-resolution}
\begin{tabular}{c||ccc}
\toprule
 & \texttt{0} & \texttt{1} & \texttt{Z} \\
\midrule
\texttt{0} & 0 & \textbf{X} & 0 \\
\texttt{1} & \textbf{X} & 1 & 1 \\
\texttt{Z} & 0 & 1 & Z \\
\bottomrule
\end{tabular}
\end{table}

\subsection*{Equations}

\begin{align*}
\text{no shift:} \quad & \texttt{Xtnd32}_i = I_i && i = 0\ldots15 \\
                       & \texttt{Xtnd32}_i = \texttt{Xsel}\cdot I_{15} && i = 16\ldots31 \\
\text{with $\times4$:} \quad & \texttt{Xtnd32}_i = E_{i-2}, \quad \texttt{Xtnd32}_1 = \texttt{Xtnd32}_0 = 0
\end{align*}

% ONE SENTENCE, required: why can the TriBuffer's behaviour NOT be written as a
% Boolean equation in AND, OR and NOT? Refer to its output range.

\section{Design}
% Answer all three:
%   (a) Why did you build ImmedXtnd in two stages (extend, then shift) rather
%       than as one expression? What does the RTL Viewer show?
%   (b) Did you write ALUsftMUX directly, or instantiate your Lab 3 Mux2 with
%       W => 5? Justify the choice. Reusing a tested component is usually the
%       better engineering answer - say whether you agree.
%   (c) What does 'Z' mean physically, and what does 'X' tell you when you see
%       it in a waveform?

\section{Implementation}

\vhdlfile{code/lab04/ImmedXtnd.vhd}{\texttt{ImmedXtnd} --- 16 to 32 bit immediate extension}

\vhdlfile{code/lab04/TriBuffer.vhd}{\texttt{TriBuffer} --- 32-bit tri-state buffer}

\section{Verification}
% Vectors, then the coverage argument.
% The argument for this lab has a specific shape and you must state it:
%   TEST WHERE THE IMPLEMENTATIONS DISAGREE, NOT WHERE THEY AGREE.
% An immediate of 0x0001 gives the same answer either way and proves nothing.

\begin{table}[htbp]
\centering
\caption{Test vectors for \texttt{ImmedXtnd} and the shared bus.}
\label{tab:lab04-vectors}
\begin{tabular}{cllll}
\toprule
\# & inputs & expected & what it proves & what it rules out \\
\midrule
1 & & & & \\
2 & & & & \\
3 & & & & \\
\bottomrule
\end{tabular}
\end{table}

\section{Results}
% Evidence: E1, E2, E7.

\figplaceholder{the waveform showing the shared bus in all three states ---
floating at \texttt{Z}, driven by one buffer, and in contention at \texttt{X}}
% \resultfig{figs/lab04/bus.png}
%   {From t = ... the bus floats at 'Z' with no driver enabled; from t = ... it
%    reads driver A; from t = ... both drivers are enabled with data differing
%    in every bit, so every bit resolves to 'X'.}
%   {fig:lab04-bus}

\figplaceholder{the RTL Viewer for \texttt{ImmedXtnd}, showing the extend stage
and the shift stage separately}

% E1: error and warning counts. State explicitly that there is no
%     inferred-latch warning.
% E4: register count for ImmedXtnd. Must be ZERO - it is combinational.

\section{Discussion}
% REQUIRED, and it must answer this question explicitly:
%
%     Why does ADDI sign-extend its immediate but LUI zero-extend its own?
%
% Answer in terms of what the bits MEAN in each instruction, then say what would
% go wrong in the Lab 12 test program if ADDI zero-extended instead.
%
% Then the usual: one bug you hit, how you found it, what it cost.
